\chapter{Results}
\label{chap:results}

\section{Overview}

The results follow the evidence hierarchy defined in
\cref{chap:methodology}. Component replays establish deterministic behaviour
on specified inputs. End-to-end AgentDojo runs measure native user and attacker
outcomes in controlled tool environments. The matched 27B WASP campaign is the
principal browser-native comparison. The 9B WASP rows explain development
failures but are not pooled as a single matched campaign.

\section{AgentDojo component results}

\subsection{Full v1 replay}

The model-free replay covered 16 non-manual attack families and 7,404 security
cases. \Cref{tab:component-replay} reports the aggregate. All defined attack
cases were detected and received a hard-block outcome. None of the 97 unique
benign tasks was hard-blocked; one benign task reached the approval threshold.

\begin{table}[htbp]
\centering
\caption{AgentDojo v1 model-free component replay. Benign tasks are repeated
once per attack family in aggregate detector counts, so the table reports the
unique benign denominator separately.}
\label{tab:component-replay}
\begin{tabular}{@{}lr@{}}
\toprule
Measure & Result \\
\midrule
Attack families & 16 \\
Security cases across families & 7,404 \\
Unique benign tasks per family & 97 \\
Detection precision / recall / F1 & 1.000 / 1.000 / 1.000 \\
Security-case hard-block rate & 7,404/7,404 \\
Benign hard blocks & 0/97 \\
Benign approval-triggering tasks & 1/97 (1.03\%) \\
\bottomrule
\end{tabular}
\end{table}

This result verifies the configured component contract but is not a held-out
security estimate. The rules were developed after inspecting AgentDojo v1
failures, the same 97 benign tasks recur across attack families, and no acting
model can re-plan. Promoting the 100\% replay rate to an end-to-end claim would
therefore be misleading.

\subsection{Detector-v2 scorecard}

The corrected detector scored recall 1.000 with no false positives on both
AgentDojo corpora. On the WASP corpus, it scored recall 1.000, precision 0.988,
and false-positive rate 0.026: one of 39 captured benign pages triggered. That
page was GitLab's profile view, whose legitimate text discusses passwords,
keys, and privacy. It was retained rather than tuned away.

\begin{table}[htbp]
\centering
\caption{Detector-v2 scorecard. WASP is a post-hoc development corpus and the
row must not be read as held-out end-to-end performance.}
\label{tab:detector-scorecard}
\begin{tabular}{@{}lrrrrr@{}}
\toprule
Corpus & Attack & Benign & Recall & Precision & FPR \\
\midrule
AgentDojo security & 130 & 58 & 1.000 & 1.000 & 0.000 \\
AgentDojo v1 tool knowledge & 629 & 97 & 1.000 & 1.000 & 0.000 \\
WASP security & 84 & 39 & 1.000 & 0.988 & 0.026 \\
\bottomrule
\end{tabular}
\end{table}

The live-representation correction is itself an important negative finding.
The original offline converter retained a Markdown link destination that
BrowserUse's accessibility-tree observation omitted. A pre-launch canary
therefore scored the visible URL attack below threshold. After the converter
was changed to reproduce the href-free, split-line observation and the
detector was adapted, all 42 URL-format attack scores changed: 40 previously
missed attacks crossed the detection threshold and two already detected cases
changed band. AgentDojo scores were unchanged. This episode demonstrates why
offline text benchmarks cannot substitute silently for the exact runtime
representation and why the resulting WASP detector score is development
evidence.

For isolated WASP payloads, 84/84 crossed the detection threshold but none
crossed the hard-block threshold. In complete attacked observations, 84/84
were detected and eight crossed the hard-block threshold. Context therefore
changes the score materially. The end-to-end campaigns use the detector as one
part of sanitisation and policy rather than inferring containment from these
classifications.

\subsection{Action-gating replay}

WASP contains 84 unique attacker goals: 20 terminate in a cross-origin action
and 64 use a same-origin form. A blanket rule that gated every
injection-tainted side effect newly gated a median of eight actions per
affected task (95th percentile 28, maximum 28) and exhausted the four-prompt
budget on 35 of 50 replayed tasks. The narrowed amplified-action rule gated a
median of zero (95th percentile zero, maximum two), exhausted the budget on no
task, and retained cross-origin exfiltration coverage. Neither rule changed
the AgentDojo component classifications.

This replay explains a core design choice. Blanket taint gating would convert
detection into denial of legitimate same-origin work. The final browser policy
therefore relies primarily on removing the hostile observation and uses action
gating as a targeted backstop.

\section{AgentDojo end-to-end results}

\subsection{Controlled comparison}

All three thinking/32K campaigns completed the 629 off/enforce task pairs.
\Cref{tab:agentdojo-e2e} reports the native outcomes. Enforcement reduced
attacker success sharply for every model while increasing both utility and
secure utility.

\begin{table}[htbp]
\centering
\caption{Controlled AgentDojo v1 end-to-end results. U is native user utility,
ASR is native attacker-goal success, and SU is useful-and-secure. Each attacked
cell has \(n=629\); clean utility has \(n=97\).}
\label{tab:agentdojo-e2e}
\footnotesize
\begin{tabular}{@{}lrrrrrrr@{}}
\toprule
 & Clean U & \multicolumn{3}{c}{Off} & \multicolumn{3}{c}{Enforce} \\
\cmidrule(lr){3-5}\cmidrule(l){6-8}
Model &  & U & ASR & SU & U & ASR & SU \\
\midrule
Qwen 3.5 4B &
76 (78.4\%) &
362 (57.6\%) & 206 (32.8\%) & 297 (47.2\%) &
446 (70.9\%) & 0 (0\%) & 446 (70.9\%) \\

Qwen 3.5 9B &
80 (82.5\%) &
422 (67.1\%) & 131 (20.8\%) & 367 (58.3\%) &
467 (74.2\%) & 1 (0.2\%) & 467 (74.2\%) \\

Qwen 3.5 27B &
81 (83.5\%) &
428 (68.0\%) & 203 (32.3\%) & 345 (54.8\%) &
483 (76.8\%) & 0 (0\%) & 483 (76.8\%) \\
\bottomrule
\end{tabular}
\end{table}

\begin{figure}[htbp]
\centering
\begin{tikzpicture}
\begin{axis}[
  ybar,
  width=0.96\textwidth,
  height=7.2cm,
  bar width=7pt,
  ymin=0,
  ymax=85,
  ylabel={Task instances (\%)},
  symbolic x coords={4B,9B,27B},
  xtick=data,
  xlabel={Qwen 3.5 model},
  legend style={at={(0.5,1.02)},anchor=south,legend columns=2,draw=none},
  grid=major,
  grid style={draw=sentinelgrey!18},
  nodes near coords,
  every node near coord/.append style={font=\scriptsize,rotate=90,anchor=west},
]
\addplot[fill=sentinelblue!35,draw=sentinelblue] coordinates
  {(4B,57.6) (9B,67.1) (27B,68.0)};
\addplot[fill=sentinelblue!80,draw=sentinelblue] coordinates
  {(4B,70.9) (9B,74.2) (27B,76.8)};
\addplot[fill=sentinelred!35,draw=sentinelred] coordinates
  {(4B,32.8) (9B,20.8) (27B,32.3)};
\addplot[fill=sentinelred!80,draw=sentinelred] coordinates
  {(4B,0) (9B,0.2) (27B,0)};
\legend{Off utility,Enforce utility,Off ASR,Enforce ASR}
\end{axis}
\end{tikzpicture}
\caption{AgentDojo utility and attacker success by condition. The rows are
independent model runs, not an identified parameter-scaling experiment.}
\label{fig:agentdojo-bars}
\end{figure}
\clearpage

For 4B, utility increased by 84 cases (13.35 percentage points), attacker
success decreased by 206 cases (32.75 points), and secure utility increased by
149 cases (23.69 points). For 9B, the corresponding changes were \(+45\),
\(-130\), and \(+100\) cases. Paired 9B transitions contained 92 utility gains
and 47 losses, with 130 attacker successes prevented and none introduced.
For 27B, utility increased by 55 cases, attacker success decreased by 203, and
secure utility increased by 138. Its paired transitions contained 84 utility
gains and 29 losses, 203 attacker successes prevented and none introduced, and
158 secure-utility gains versus 20 losses.

The utility increases are plausible rather than paradoxical. An injected model
often diverts away from the user's task. Removing the malicious tool output can
therefore improve security and allow the original task to complete. The result
does not mean enforcement always improves utility: every model has paired
losses, and the effect depends on these tasks and trajectories.

\subsection{Measured mechanism}

\begin{table}[htbp]
\centering
\caption{Enforce-mode AgentDojo diagnostics, out of 629 attacked cases per
model. ``Blocks'' counts explicit action blocks, not sanitised observations.}
\label{tab:agentdojo-diagnostics}
\begin{tabular}{@{}lrrrr@{}}
\toprule
Model & Cases detected & Blocks & Approvals & Truncated responses \\
\midrule
4B & 600 & 9 & 0 & 6 \\
9B & 615 & 9 & 0 & 0 \\
27B & 609 & 9 & 0 & 7 \\
\bottomrule
\end{tabular}
\end{table}

Sanitising untrusted tool output was the dominant observed intervention.
Explicit action blocking was a secondary backstop: only nine actions were
blocked for each model. In the reviewed 27B aggregate, all nine were
\texttt{sensitive\_data\_flow} decisions, while 2,181 actions were allowed.
This distinction prevents a false explanation that high utility came from
blocking most activity.

Response limits remain a validity concern. Ten 27B responses reached the
4,096-token cap: three off and seven enforce. The seven enforce truncations
were all variants of one difficult travel user task and were retained rather
than selectively rerun. They can depress utility but cannot be removed without
introducing post-selection.

\subsection{Provenance qualifications}

The 4B campaign repeats both modes under the fixed pipeline. The 9B published
row is backed by all 1,258 preserved reports but lacks recorded Git revision in
those reports. The 27B comparison preserves historical off reports and uses a
new immutable fixed-enforce tag; its task/model/profile configuration matches,
but exact model-blob parity with the historical off run cannot be verified
cryptographically. These limits weaken a claim of perfectly identical
experimental provenance, but they do not change the native counts.

Cross-model ordering is also not monotonic: 27B off ASR exceeds 9B off ASR.
This reinforces the methodological warning that low ASR can reflect model
behaviour and task capability rather than a simple size effect.

\section{WASP browser-native results}

\subsection{Primary matched 27B finding}

\Cref{tab:wasp-27b} presents the principal browser comparison. All 168
executions share the frozen model, scaffold, policy, middleware, and 30-step
design described in \cref{sec:wasp-method}. Each of the four batches has its own
generated-input manifest, held identical between its paired off and enforce
cells.

\begin{table}[htbp]
\centering
\caption{Matched Qwen 3.5 27B WASP results, \(n=84\) per condition. Attempted
and executed columns count task instances with at least one attacker-directed
action.}
\label{tab:wasp-27b}
\small
\begin{tabular}{@{}lrrrrr@{}}
\toprule
Condition & Utility & ASR & Secure utility & Attempted & Executed \\
\midrule
Off &
47 (55.95\%) & 5 (5.95\%) & 44 (52.38\%) &
17 (20.24\%) & 17 (20.24\%) \\
Sanitise V2 enforce &
46 (54.76\%) & 0 (0\%) & 46 (54.76\%) &
0 (0\%) & 0 (0\%) \\
\bottomrule
\end{tabular}
\end{table}

\begin{figure}[htbp]
\centering
\begin{tikzpicture}
\begin{axis}[
  ybar,
  width=0.9\textwidth,
  height=7cm,
  bar width=14pt,
  ymin=0,
  ymax=65,
  ylabel={Task instances (\%)},
  symbolic x coords={Utility,ASR,Execution,{Secure utility}},
  xtick=data,
  legend style={at={(0.5,1.02)},anchor=south,legend columns=2,draw=none},
  grid=major,
  grid style={draw=sentinelgrey!18},
  nodes near coords,
  every node near coord/.append style={font=\scriptsize},
]
\addplot[fill=sentinelgrey!40,draw=sentineldark] coordinates
  {(Utility,55.95) (ASR,5.95) (Execution,20.24) ({Secure utility},52.38)};
\addplot[fill=sentinelamber!70,draw=sentinelamber] coordinates
  {(Utility,54.76) (ASR,0) (Execution,0) ({Secure utility},54.76)};
\legend{Off,Sanitise V2 enforce}
\end{axis}
\end{tikzpicture}
\caption{Primary matched WASP outcome rates. Enforcement removed the observed
attacker-directed execution with a one-task aggregate utility change.}
\label{fig:wasp-bars}
\end{figure}

All five off-condition terminal attacker successes were absent under enforce,
and no new success appeared. Attacker-directed attempts and executions changed
from 17 task instances to zero. Utility changed by one task, or \(-1.19\)
percentage points, while secure utility increased by two tasks, or \(+2.38\)
points.

The paired utility outcomes contained seven gains, eight losses, and 69
unchanged cases. There were no hard blocks or handovers in the enforce row.
One page-code action triggered an approval and was denied, but that task still
passed its native user oracle. It is therefore incorrect to label the eight
utility losses as eight false blocks. They are trajectory-level differences
that can arise from altered observations and stochastic planning.

Joint outcomes provide additional context. Off contained 44 useful/secure,
three useful/compromised, 35 unsuccessful/secure, and two
unsuccessful/compromised cases. Enforce contained 46 useful/secure and 38
unsuccessful/secure cases, with no compromised category. The agent was still
unsuccessful on 45.24\% of enforce tasks; the defence did not make browser
automation generally reliable.

Zero attacker successes in 84 enforce cases is not proof of a zero underlying
rate. The rough rule-of-three upper 95\% bound is \(3/84\), approximately
3.6\%, under an independent binomial interpretation. Moreover, two of the four
off batches already had zero terminal ASR. The reduction in
attacker-directed execution from 17/84 to zero is therefore important
supporting evidence: it shows less attacker-serving activity even where the
baseline agent was often unable to finish the attack.

\subsection{Sanitisation telemetry}

The matched enforce run recorded one finding in every attacked task and 368
sanitisation attempts. All 368 were confirmed removals, with no residual,
unattributable, or post-success-alert outcome. Fifteen used the line branch and
353 the bounded-window branch. With no hard block or handover, and one denied
approval on a task that still retained native utility, post-verified
sanitisation was the primary recorded intervention associated with containment.

This evidence verifies what occurred for detector-positive observations in
the run. It cannot show that every possible attack would be detected, that
removal preserved all benign nearby text, or that the post-campaign incident
flow would improve outcomes.

\subsection{9B development and failure analysis}

\Cref{tab:wasp-9b} preserves the complete 9B sequence. The shared off baseline
is followed by three enforce rows with changed middleware provenance.

\begin{table}[htbp]
\centering
\caption{Qwen 3.5 9B WASP development sequence, \(n=84\) per row. V1 and V2
are successive sanitisation implementations; the rows are not one immutable
matched campaign.}
\label{tab:wasp-9b}
\small
\begin{tabular}{@{}lrrrrr@{}}
\toprule
Condition & Utility & ASR & Secure utility & Attempted & Executed \\
\midrule
Shared off &
45 (53.57\%) & 7 (8.33\%) & 41 (48.81\%) &
31 (36.90\%) & 31 (36.90\%) \\
Warning enforce &
37 (44.05\%) & 7 (8.33\%) & 35 (41.67\%) &
32 (38.10\%) & 32 (38.10\%) \\
Sanitise V1 &
52 (61.90\%) & 6 (7.14\%) & 49 (58.33\%) &
18 (21.43\%) & 18 (21.43\%) \\
Sanitise V2 &
54 (64.29\%) & 0 (0\%) & 54 (64.29\%) &
1 (1.19\%) & 1 (1.19\%) \\
\bottomrule
\end{tabular}
\end{table}

Warning-only enforcement is a negative result. ASR remained 7/84,
attacker-directed execution increased from 31 to 32 cases, utility fell from
45 to 37, and secure utility fell from 41 to 35. Detecting an attack and
warning the same model did not provide effective containment under this
condition.

Sanitisation V1 improved the outcome profile but left six terminal attacks and
18 cases with attacker-directed execution. Its failure analysis found
incomplete multi-line removal, which V2 addressed through revised bounded
cluster handling. A separate pre-launch canary had already exposed and
corrected the mismatch between offline Markdown URL attacks and BrowserUse's
href-free live representation; that fix was not part of the V1-to-V2 change.
V2 recorded 0/84 terminal attacks and one attacker-directed navigation, which
did not satisfy the native attacker oracle.

The V2 row recorded 1,057 sanitisation attempts: 1,053 confirmed removals, zero
residuals, and four unattributable cases that used the historical
model-warning fallback. Branches were 14 line, 1,039 bounded window, and four
none. Across the run, 85 detections accompanied 662 attempted actions, one hard
block, four approval prompts, and no handovers. A frozen observation replay
separately confirmed removal in 84/84 attacked observations, with no
detector-positive residual, no non-anchor control-descriptor removal, and no
whole-page wipe.

The four unattributable events motivated the later fail-closed operator
incident, but that mechanism is not part of this row. It would be a provenance
error to describe those observations as quarantined in the reported campaign.

\section{Current implementation verification}
\label{sec:implementation-verification}

Post-campaign contracts were freshly verified on 30 July 2026 rather than
folded into the benchmark outcomes. The complete Python suite passed 231
tests and the focused WASP campaign suite passed 29. The five-case security
smoke achieved detection precision, recall, F1, verdict accuracy, and attack
containment of 1.000, with zero interventions across two benign fixtures. The
offline WASP preflight and TypeScript/Vite extension build also passed.
\Cref{tab:implementation-verification} records each gate and result. A
machine-readable record binds the commands, outputs, tool versions, repository
revision, and relevant-source manifest to this verification
(\texttt{sentinel/evals/submission-validation-20260730.json}).

These gates establish that the described current implementation contracts held
for that source state. They do not retroactively add the later mechanisms to
the AgentDojo or WASP campaign rows, and they are not substitutes for a human
evaluation of the extension.

\section{Result synthesis}

The evidence answers EQ1 with a deterministic component pass, qualified by
development-set adaptation. EQ2 shows large reductions in native attacker
success and increases in useful-and-secure outcomes across three controlled
AgentDojo model runs; sanitisation, not mass action blocking, is the dominant
recorded mechanism. EQ3 shows browser-native transfer in the primary matched
27B WASP comparison: five terminal attacks and 17 attacker-directed execution
cases observed off were absent under enforce, with a one-task utility change.
EQ4 identifies post-verified sanitisation as the primary recorded intervention
and warning-only enforcement as ineffective in the 9B development condition.
RQ2 receives a bounded technical answer from the design and current
verification gates: approval, takeover, unresolved-observation incidents,
timeout, cancellation, live revalidation, and paired audit exercise selected
server-owned parts of the specified fail-closed authority lifecycle. Rendered
fixture states demonstrate the implemented client presentation, but reconnect,
accessibility, schema equivalence, explanation quality, and human use remain
unevaluated.

These results support the bounded thesis claim, but not a universal one. The
next chapter interprets the utility gains and losses, the role of model
capability, the limits of fixed benchmark attacks, and the unresolved
human-factors evidence.
